\section{Transformasi Ekuivalensi Bersyarat (Implikasi \& Biimplikasi)}

Ekspresi logika sering menggunakan operator implikasi ($\rightarrow$) dan biimplikasi ($\leftrightarrow$). Untuk menyederhanakannya, kedua bentuk tersebut dapat ditransformasikan ke operator dasar AND, OR, dan NOT \cite{rosen2019}.

\subsection{Transformasi Implikasi}

Implikasi $(p \rightarrow q)$ merepresentasikan pernyataan ``Jika $p$, maka $q$''. Bentuk ini mempunyai ekuivalensi standar berikut:

\begin{teorema}[Ekuivalensi Implikasi]
Struktur $(p \rightarrow q)$ ekuivalen dengan bentuk disjungsi $(\neg p \vee q)$.
\end{teorema}

Bentuk yang sama juga dapat dituliskan sebagai:
\[
p \rightarrow q \equiv \neg (p \wedge \neg q)
\]

\textbf{Contoh singkat:} \\
``Jika Anda menyumbang darah ($p$), maka saya memberi apresiasi ($q$).'' \\
Secara logika, pernyataan tersebut ekuivalen dengan ``Anda tidak menyumbang darah ($\neg p$), atau saya memberi apresiasi ($q$)'', yaitu $\neg p \vee q$.

\subsection{Transformasi Biimplikasi}

Biimplikasi menyatakan hubungan dua arah: ``jika dan hanya jika''.

\begin{teorema}[Ekuivalensi Biimplikasi]
Struktur $(p \leftrightarrow q)$ dapat diuraikan sebagai:
\begin{align*}
    p \leftrightarrow q &\equiv (p \rightarrow q) \wedge (q \rightarrow p) \\
    \text{dengan substitusi implikasi: } & \\
    p \leftrightarrow q &\equiv (\neg p \vee q) \wedge (\neg q \vee p)
\end{align*}
\end{teorema}
Bentuk ekuivalen lain yang lazim digunakan adalah:
\[
p \leftrightarrow q \equiv (p \wedge q) \vee (\neg p \wedge \neg q)
\]

\subsection{Pola Penyederhanaan Campuran}

Dalam praktik, terdapat bentuk campuran yang sering muncul dan dapat disederhanakan secara cepat:

\begin{align*}
    p \vee (\neg p \wedge q) &\equiv (p \vee q) \\
    p \wedge (\neg p \vee q) &\equiv (p \wedge q)
\end{align*}

Pembuktian untuk bentuk pertama:
\begin{enumerate}
    \item Terapkan distributif: $\equiv (p \vee \neg p) \wedge (p \vee q)$
    \item Terapkan hukum negasi: $\equiv \mathbf{T} \wedge (p \vee q)$
    \item Terapkan identitas: $\therefore \equiv (p \vee q)$
\end{enumerate}
Pola-pola ini bermanfaat untuk mempercepat penyederhanaan ekspresi pada analisis logika maupun refactoring kode.
